\scanpage{280}
- \mspage{33} -
\margin{IX}
de traits \(S'\longrightarrow S.\) Enfin, si (iv) est satisfaite pour \(A_K,\) elle l'est aussi manifestement pour tout sous-schéma abélien et tout schéma abélien quotient de \(A_K,\) et en particulier pour tout schéma abélien \uline{isogène} à \(A\) (par exemple pour \(A'_K\)).

\uline{Corollaire} 3.5.2. \uline{Soit} \(g\) \uline{un générateur topologique de} \(\mathbb Z_\ell(1)=T_\ell(\bar k^*).\) \uline{Les conditions équivalentes de} 3.5 \uline{équivalent aussi à la condition suivante} :

(v) \uline{L'action du groupe d'inertie} \(I\) \uline{sur} \(U=T_\ell(A(\bar K))\) \uline{se fait à travers le plus grand quotient pro-\(\ell\)-fini} \(I(\ell)\) \uline{de} \(I,\) \uline{qui est canoniquement isomorphe à} \(\mathbb Z_\ell(1)\) (I\qquad), \uline{et l'opération} \(g_U\) \uline{du générateur} \(g\) \uline{sur} \(U\) \uline{satisfait la condition}
\[
(3.5.3)\qquad(1-g_U)^2=0.
\]
En effet, pour toute action unipotente continue d'un groupe profini \(I\) sur un espace vectoriel \(\ell\)-adique, on voit aussitôt que l'image de \(I\) est un pro-\(\ell\)-groupe. Il s'ensuit que (v) équivaut à (iv).

Nous pouvons maintenant prouver le résultat central du présent exposé :

\uline{Théorème} 3.6. (\uline{théorème de réduction semi-stable}). \uline{Soient} \(S\) \uline{un schéma noethérien régulier connexe de dimension} 1, \(K\) \uline{le corps des fonctions rationnelles de} \(S,\) \(A_K\) \uline{un schéma abélien sur} \(K.\) \uline{Alors il existe une extension finie galoisienne} \(K'\) \uline{de} \(K,\) \uline{telle que le schéma abélien} \(A_{K'}=A_K\otimes_K K'\) \uline{ait réduction semi-stable} (3.4) \uline{sur le normalisé} \(S'\) \uline{de} \(S\) \uline{dans} \(K'.\)
\figure{\begin{tikzpicture}\draw (0,0)--(.2,0);\draw (2,.5)--(2.2,.5) (2.05,.45) arc (180:360:.05);\end{tikzpicture}}
\margin{351}
\margin{280}
\margin{351}
